% Generado por sistema/backend ReporteTex.java -- no editar a mano.
% Datos: Base de datos: 148 importado (titulados) / Base de datos: 53 importado (empleadores)
\subsubsection*{C1. Situación laboral $\times$ satisfacción con la formación}

\textit{Propósito.} Determinar si la valoración de la formación recibida cambia según la situación laboral actual del titulado, como evidencia del resultado del proceso formativo frente a las condiciones de empleo. Respalda el criterio ARCU-SUR 3.2.1 / 3.2.3. Se aplica sobre: todos los titulados, con nivel de significancia $\alpha = 0.05$.

\textit{Paso 1. Contar con la muestra: tabla de contingencia} (frecuencia observada $FO_{ij}$, $N = 147$).

\begin{mitabla}{C1: frecuencia observada $FO_{ij}$}
\begin{tabular}{lcc|c}
\toprule
Situación laboral & En desacuerdo & De acuerdo & Total fila \\
\midrule
Trabaja en una organización & 27 & 81 & 108 \\
Emprendimiento propio & 5 & 9 & 14 \\
No trabaja & 6 & 19 & 25 \\
\midrule
Total columna & 38 & 109 & $N=147$ \\
\bottomrule
\end{tabular}
\end{mitabla}

\textit{Paso 2. Calcular la frecuencia esperada} de cada celda, $FE_{ij} = (\text{total fila } i)(\text{total columna } j)/N$. Por ejemplo, para la celda ``Trabaja en una organización -- En desacuerdo'': $FE_{11} = 108 \times 38 / 147 = 27.9184$.

\begin{mitabla}{C1: frecuencia esperada $FE_{ij}$}
\begin{tabular}{lcc}
\toprule
Situación laboral & En desacuerdo & De acuerdo \\
\midrule
Trabaja en una organización & $108\cdot38/147 = 27.92$ & $108\cdot109/147 = 80.08$ \\
Emprendimiento propio & $14\cdot38/147 = 3.62$ & $14\cdot109/147 = 10.38$ \\
No trabaja & $25\cdot38/147 = 6.46$ & $25\cdot109/147 = 18.54$ \\
\bottomrule
\end{tabular}
\end{mitabla}

\textit{Paso 3. Calcular el estadístico chi-cuadrado,} sumando $(FO_{ij}-FE_{ij})^2/FE_{ij}$ de cada celda.

\begin{mitabla}{C1: cálculo de $X_0^2$ celda por celda}
\begin{tabular}{L{5.4cm}cccc}
\toprule
Celda & $FO_{ij}$ & $FE_{ij}$ & $(FO_{ij}-FE_{ij})^2$ & $(FO_{ij}-FE_{ij})^2/FE_{ij}$ \\
\midrule
Trabaja en una organización -- En desacuerdo & 27 & 27.9184 & 0.8434 & 0.0302 \\
Trabaja en una organización -- De acuerdo & 81 & 80.0816 & 0.8434 & 0.0105 \\
Emprendimiento propio -- En desacuerdo & 5 & 3.6190 & 1.9070 & 0.5269 \\
Emprendimiento propio -- De acuerdo & 9 & 10.3810 & 1.9070 & 0.1837 \\
No trabaja -- En desacuerdo & 6 & 6.4626 & 0.2140 & 0.0331 \\
No trabaja -- De acuerdo & 19 & 18.5374 & 0.2140 & 0.0115 \\
\midrule
\multicolumn{4}{r}{$X_0^2 = \sum$} & 0.7960 \\
\bottomrule
\end{tabular}
\end{mitabla}

\textit{Paso 4. Determinar el estadístico teórico.} Con $k = 3$ filas y $m = 2$ columnas, los grados de libertad son $(k-1)(m-1) = (3-1)(2-1) = 2$. De la tabla de chi-cuadrada: $X_{0.05,\,2}^2 = 5.991$.

\textit{Paso 5. Establecer la relación entre estadísticos:} $X_0^2 < X_{0.05,\,2}^2 \;\Rightarrow\; 0.7960 < 5.991$.

\textit{Paso 6. Conclusión.}

\begin{center}
\begin{tabular}{l}
Si ($X_0^2 < X_{0.05,\,2}^2$)\\
\quad \fbox{Independencia}\\
Sino\\
\quad Relación entre las variables\\
FinSi
\end{tabular}
\end{center}

No se puede rechazar que situación laboral y satisfacción sean independientes.

\textit{Observación:} 1 de 6 celdas tienen $FE < 5$, por lo que la prueba pierde confiabilidad y la conclusión debe tomarse con cautela.

